\subsection{Inverse-Problem-Like Formulations in Chest Pathology Imaging}

While chest X-rays are not inverse problems in the strict physical sense (as reconstruction from projections is handled by the scanner), 
several problem formulations in chest pathology imaging exhibit inverse characteristics where 
the underlying cause (disease, lesion, or anatomy) must be inferred from incomplete or noisy effects (observed image data). 
Below we summarize key examples.

\paragraph{Denoising and Artifact Removal.}
Given a noisy or artifact-corrupted chest X-ray \( y = x + \epsilon \), 
the goal is to recover the clean image \( x \). 
A diffusion prior acts as a strong denoiser:
\[
x^* = \arg\max_x \, p(y|x)\,p_\theta(x),
\]
where \( p_\theta(x) \) constrains the solution to anatomically plausible radiographs. 
\textit{Application:} low-dose or mobile X-ray correction.

\paragraph{Limited-Angle or Sparse-View Reconstruction.}
When only partial X-ray projections are available (e.g., for dose reduction or portable imaging),
the forward model is:
\[
y = A(x) + \epsilon,
\]
with \( A \) representing the limited projection operator. 
The inverse task is to reconstruct the full radiograph \( x \) consistent with \( y \), 
regularized by a diffusion prior \( p_\theta(x) \). 
\textit{Application:} dose-efficient imaging and fast screening systems.

\paragraph{2D-to-3D Chest Reconstruction.}
Given a 2D projection \( y \) of an unknown 3D volume \( x \), 
the physical process can be described as:
\[
y = A_\text{proj}(x),
\]
where \( A_\text{proj} \) denotes the X-ray projection operator. 
The goal is to infer \( x \) that could plausibly produce \( y \). 
A generative model supplies the anatomical prior \( p_\theta(x) \) over valid 3D chest geometries. 
\textit{Application:} reconstructing volumetric lung structure or estimating lesion depth from a single X-ray.

\paragraph{Semantic Inference of Pathologies.}
Here, the observed image \( y \) is used to infer the latent pathology state \( z \), 
with a generative decoder \( p_\theta(y|z) \) modeling disease appearance:
\[
p(z|y) \propto p(y|z)p(z).
\]
This constitutes an \textit{inverse inference} problem in latent space, 
where the goal is to infer the hidden pathological cause from observed image evidence. 
\textit{Application:} disease localization, severity estimation, and progression modeling.

\paragraph{Compositional Generative Inference (Our Formulation).}
In this research, we introduce a compositional generative inference framework for chest X-rays, 
where pretrained diffusion models trained on disjoint datasets (e.g., TB vs.~Normal) are combined at inference time. 
Formally, let \( s_\text{TB}(x,t) \) and \( s_\text{Normal}(x,t) \) denote the score fields 
from the respective pretrained models. 
Following the \emph{SuperDiff} formulation, a composite score field is defined as:
\[
s_\text{comp}(x,t) = \kappa\, s_\text{TB}(x,t) + (1 - \kappa)\, s_\text{Normal}(x,t),
\]
where \( \kappa \in [0,1] \) controls the compositional weighting. 
The corresponding reverse diffusion process integrates this composite score to synthesize novel samples that 
lie on the shared manifold between the TB and Normal data distributions. 

\textit{Application:} 
generation of intermediate or novel pathology states for data augmentation, interpretability, and 
latent-space exploration of disease semantics. 
Although not an inverse problem in the classical physics sense, this formulation constitutes an 
\emph{inverse-problem-like inference} in the latent semantic space of disease composition, 
where the goal is to infer plausible intermediate states from disjoint learned priors.